\subsection{FUNCTORIAL PROPERTIES}

In this no., all algebras are assumed to be associative and unital and every algebra homomorphism is assumed to be unital.

PROPOSITION 9. Let A, B be two graded K-algebras, E an (A, A)-bimodule and F a graded (B, B)-bimodule; let \(\rho: \mathbf{A} \to \mathbf{B}\) be a graded algebra homomorphism and \(\theta: \mathbf{E} \to \mathbf{F}\) a graded A-homomorphism of A-bimodules (relative to \(\rho\) ), of degree 0. Then:

\begin{enumerate}
\def\labelenumi{(\roman{enumi})}
\item
  For every \(\varepsilon\) -derivation \(d':\mathbf{B}\to\mathbf{F}\) , \(d'\circ\rho:\mathbf{A}\to\rho^{*}(\mathbf{F})\) is an \(\varepsilon\) -derivation of the same degree.
\item
  For every \(\varepsilon\) -derivation \(d: \mathbf{A} \to \mathbf{E}\) , \(\theta \circ d: \mathbf{A} \to \rho^{*}(\mathbf{F})\) is an \(\varepsilon\) -derivation of the the same degree.
\end{enumerate}

The two assertions follow immediately from the relations

\[
\begin{array}{l} d ^ {\prime} (\rho (x y)) = d ^ {\prime} (\rho (x) \rho (y)) = d ^ {\prime} (\rho (x)) \rho (y) + \varepsilon_ {\delta ^ {\prime}, \xi} \rho (x) d ^ {\prime} (\rho (y)) \\ \theta (d (x y)) = \theta ((d x) y + \varepsilon_ {\delta, \xi} x (d y)) = \theta (d x) \rho (y) + \varepsilon_ {\delta, \xi} \rho (x) \theta (d y) \end{array}
\]

for \(x \in \mathbf{A}\) homogeneous of degree \(\xi\) and \(y \in \mathbf{A}\) , \(\delta\) and \(\delta'\) denoting the respective degrees of \(d\) and \(d'\) .

COROLLARY. Let S be a generating system of the algebra A. In order that \(d' \circ \rho = \theta \circ d\) , it is necessary and sufficient that \(d'(\rho(x)) = \theta(d(x))\) for all \(x \in S\) .

This is an immediate consequence of Proposition 9 and no. 5, Corollary to Proposition 3.

Under the conditions of Proposition 9, we know that B has (by means of \(\rho\) ) an (A, A)-bimodule structure (II, § 1, no. 14, Example 1).

PROPOSITION 10. Under the conditions of Proposition 9, for an \(\varepsilon\) -derivation \(d': \mathbf{B} \to \mathbf{F}\) to be A-linear for the left (resp. right) A-module structures on B and \(\rho_*(\mathbf{F})\) , it is necessary and sufficient that \(d'\) be zero on the subalgebra \(\rho(\mathbf{A})\) of B.

We perform the proof for left A-module structures. For \(a \in \mathbf{A}\) , \(b \in \mathbf{B}\) ,

\[
d ^ {\prime} (\rho (a) b) = d ^ {\prime} (\rho (a)) b + \rho (a) d ^ {\prime} b
\]

and hence if \(d' \circ \rho = 0\) , \(d'\) is linear for the left A-module structures on B and \(\rho^*(\mathbf{F})\) . Conversely, if this is so, in particular

\[
d ^ {\prime} (\rho (a)) = d ^ {\prime} (\rho (a). 1) = \rho (a) d ^ {\prime} (1) = 0
\]

(no. 5, Proposition 3).

In particular let \(D_{K}(B,F)\) denote the K-module of derivations of B into F (no. 2); those among these derivations which are A-linear, in other words those which are zero on \(\rho(A)\) , form a sub-K-module of \(D_{K}(B,F)\) , denoted by \(D_{A,\rho}(B,F)\) or simply \(D_{A}(B,F)\) (obviously \(D_{K}(B,F)=D_{K,\phi}(B,F)\) , where \(\phi:K\to B\) is the homomorphism defining the K-algebra structure on B).

Now let A, B, C be three graded K-algebras, \(\rho:A\to B\) , \(\sigma:B\to C\) two graded algebra homomorphisms and G a graded (C, C)-bimodule; if \(D_{A}(B,G)\) , \(D_{B}(C,G)\) and \(D_{A}(C,G)\) denote the respective K-modules \(D_{A,\rho}(B,\sigma_{*}(G))\) , \(D_{B,\sigma}(C,G)\) and \(D_{A,\sigma_{o}\rho}(C,G)\) , \(D_{B}(C,G)\) is clearly a sub-K-module of \(D_{A}(C,G)\) since \(\sigma(\rho(A))\subset\sigma(B)\) .

PROPOSITION 11. Under the above conditions, there is an exact sequence of K-homomorphisms

\[
0 \rightarrow \mathrm{D} _ {\mathrm{B}} (\mathrm{C}, \mathrm{G}) \xrightarrow {u} \mathrm{D} _ {\mathrm{A}} (\mathrm{C}, \mathrm{G}) \xrightarrow {v} \mathrm{D} _ {\mathrm{A}} (\mathrm{B}, \mathrm{G})\tag{27}
\]

where u is the canonical injection and v the homomorphism \(d \mapsto d \circ \sigma\) (Proposition 9). The kernel of v is the set of derivations \(d: C \to G\) such that \(d(\sigma(b)) = 0\) for all \(b \in B\) , which is precisely the image of u.

